\documentclass[10pt,a4paper]{amsart}
\usepackage[margin=19mm]{geometry}
\usepackage{amsmath,amssymb,mathtools}
\usepackage[scr=rsfs,cal=euler]{mathalfa}
\usepackage{xfrac}
\usepackage[unicode,hidelinks,bookmarks=false]{hyperref}
\newcommand{\ie}{i.e.\ }
\newcommand{\cf}{cf.\ }
\newcommand{\resp}{respectively\ }
\newcommand{\Subsection}{\subsection}
\providecommand{\nobreakdash}{\nobreak\hbox{-}}
\setlength{\emergencystretch}{2em}
\multlinegap0pt
\usepackage{amssymb}
\usepackage{tikz}
\newtheorem{proposition}{Proposition}
\title{{Symmetries of Random Partitions}}
\author{Alexander Gnedin}
\date{Source: arXiv:2606.22533v1 (CC BY 4.0)}
\begin{document}
\maketitle

\noindent\emph{Source and licence.} {Symmetries of Random Partitions}. Version arXiv:2606.22533v1 (CC BY 4.0), distributed by arXiv under the Creative Commons Attribution 4.0 International licence, http://creativecommons.org/licenses/by/4.0/, as recorded in the arXiv licence metadata for this exact version and shown on the abstract page https://arxiv.org/abs/2606.22533v1. Full licence notice: \texttt{licenses/arxiv-2606.22533-CC-BY-4.0.txt}.\par\smallskip

\noindent\textit{Editorial scope.} Counted unit: the article's proposition n=4 (source0.tex lines 980--982). The article's complete original proof of it is delivered with it (source0.tex lines 983--1050, one \texttt{proof} environment of 68 lines). No mathematics is written, repaired or re-derived: the statement and the proof are the article's own lines, in the article's own notation, and no retained line is substituted or re-flowed.  Retained uncounted as the source-local context the proof needs: lines 775--781.  Nothing was omitted from any retained span, so the package carries no omission record.  No retained line contains a citation, so the wrapper carries no bibliography and no bibliography entry.  No retained line contains a figure environment, an image-inclusion command or drawing code, so no image is delivered and the package declares no asset dependency.  Editorial formatting only, in addition: an AMS \texttt{amsart} preamble that restates the article's own declarations and macros used by the retained lines, namely the environments \texttt{proposition} on separate counters, together with the standard packages those lines need.  The retained environments print in this document as Proposition 1 in order of appearance, whereas the article prints the same items with the numbering of its own full document.  The complete native source of the article as distributed in the arXiv e-print for this exact version is bundled unchanged as \texttt{sources/203354/source0.tex}.
\begin{equation}\label{SprRule}
\sum_{\pi_n:\,D_i\pi_n=\pi_{n-1}} P(\pi_n)
=
\sum_{\pi_n:\,D_n\pi_n=\pi_{n-1}} P(\pi_n),
\qquad
\pi_{n-1}\in{\mathcal P}_{n-1},\;\;1\le i\le n-1 .
\end{equation}

\begin{proposition}\label{n=4}
Every spreadable random partition $\Pi_4$ of $[4]$ is exchangeable.
\end{proposition}

\begin{proof}
Write $P(\pi):=\mathbb P[\Pi_4=\pi]$ for $\pi\in\mathcal P_4$. We verify
that the system of linear equations defined by \eqref{SprRule} implies that
$P$ is constant on each orbit ${\mathcal O}_\lambda$
of the symmetric group $\mathfrak S_4$.


{\bf 1. $\lambda=(3,1)$.} Consider $\pi_3=123\in\mathcal P_3$. The spreadability
condition \eqref{SprRule} requires that the total  mass projecting from $\mathcal P_4$
to $\pi_3$ is the same for each $D_i$. The partition $1234$ 
maps to $123$ under every $D_i$, contributing an identical $P(1234)$ term which cancels 
out. Looking at the remaining preimages of shape $(3,1)$, 
under $D_4$ the unique preimage is $123|4$. Under $D_3$, ball $3$
is deleted and ball $4$ is relabelled as $3$, so the unique preimage is
$124|3$. Likewise, under $D_2$ and $D_1$, the unique preimages are $134|2$
and $234|1$, respectively. Therefore,  \eqref{SprRule} gives
\begin{equation}\label{eq:shape31}
P(123|4)
=
P(124|3)
=
P(134|2)
=
P(1|234).
\end{equation}




{\bf 2.  $\lambda=(2,1,1)$.}
Consider  $\pi_3=1|2|3$. Applying
\eqref{SprRule} and comparing the four deletion operators yields a linear
system for the six partitions of shape $(2,1,1)$. From this we find 
\[
P(12|3|4)
=
P(13|2|4)
=
P(14|2|3)
=
P(1|23|4)
=
P(1|24|3)
=
P(1|2|34).
\]


{\bf 3. $\lambda=(2,2)$.}
Now consider the target partition $\pi_3=12|3\in\mathcal P_3$. Under $D_4$,
the preimages are $124|3$, $12|34$ and $12|3|4$. Under $D_2$, the preimages are
$123|4$, $13|24$ and $13|2|4$. Thus, using \eqref{eq:shape31} and the equalities above, \eqref{SprRule} gives
$P(12|34)=P(13|24).$
Repeating the same argument with $D_1$ we
obtain
\[
P(12|34)
=
P(13|24)
=
P(14|23).
\]

{\bf 4. $\lambda=(1,1,1,1)$ and $\lambda=(4)$.} These are the  ranked shapes corresponding to single-element orbits.


We see that spreadable $P$ must be constant on every orbit, hence $\Pi_4$ is exchangeable.
\end{proof}

\end{document}
